\begin{abstract}
Tool-using agents act through servers they do not control. In the Model
Context Protocol (MCP), a user approves a tool call by reading the server's
declaration, but the effect of the call comes from the server's implementation,
and the only evidence the client receives is a response written by that same
server. A server that is compromised after approval can write different bytes,
or write them to a different place, and still return exactly the response an
honest server would. We first show that inspecting responses cannot close this
gap. If an honest and a malicious server produce the same transcript, any
monitor that sees only the transcript flags the malicious server exactly as
often as it wrongly flags the honest one. A measurement over 8,692 public
registry servers agrees: no auditing configuration we tested reached our
pre-registered target of 20\% detection at no more than 10\% false positives.

We then present MCPGate, a client-side admission layer. The unmodified server
runs against a private copy of the state. After every writer has stopped, a
trusted mediator reads the complete result once, compares it with a contract
bound to the approved call, and either promotes exactly the bytes it checked or
discards them. We evaluate MCPGate on seven third-party servers that were
frozen before any outcome was observed, covering filesystem and SQLite effects
in both Node and Python, against an adversary that also falsifies its own
read-backs. MCPGate prevented all 41 applicable attacks and blocked no honest
workflow. A task-specific static permission policy prevented 20. No defense, a
plain sandbox, and a response auditor each prevented only the 9 attacks that
had no effect on these servers to begin with. Per-call cost under MCPGate was
indistinguishable from running without a defense. MCPGate guarantees what
enters trusted state; it does not stop effects outside the boundary it
mediates.
\end{abstract}
